\documentclass[12pt]{article}

\usepackage[utf8]{inputenc}
\usepackage[margin=1in]{geometry}
\usepackage{amsmath,amssymb,amsfonts}
\usepackage{physics}
\usepackage{graphicx}
\usepackage{hyperref}
\usepackage{cite}
\usepackage{booktabs}

\title{Area Laws in Quantum Nash Equilibria}
\author{C. L}
\date{\today}

\begin{document}

\maketitle

\begin{abstract}
Abstract here.
\end{abstract}

\section{Setup}
\label{sec:setup}

\subsection{Quantum game formulation}

We consider an $L$-player quantum game where each player $i$ controls a single qubit.
The joint state of all players is described by a quantum state $\ket{\psi} \in (\mathbb{C}^2)^{\otimes L}$.
Each player's strategy consists of a single-qubit unitary $U_i \in \mathrm{SU}(2)$ applied to their qubit.
Given a strategy profile $(U_1, \ldots, U_L)$ applied to a fiducial state $\ket{\psi_0}$, the resulting state is
\begin{equation}
    \ket{\psi} = (U_1 \otimes U_2 \otimes \cdots \otimes U_L) \ket{\psi_0}.
\end{equation}

The payoff for player $i$ is defined as the expectation value of a Hermitian operator $H_i$:
\begin{equation}
    \label{eq:payoff}
    E_i = \bra{\psi} H_i \ket{\psi}.
\end{equation}
In the quantum Prisoner's Dilemma, the Hamiltonians $H_i$ are diagonal in the computational basis,
\begin{equation}
    H_i = \sum_{\mathbf{s} \in \{0,1\}^L} h_i(\mathbf{s}) \ketbra{\mathbf{s}}{\mathbf{s}},
\end{equation}
where $h_i(\mathbf{s})$ assigns a real-valued payoff to each measurement outcome $\mathbf{s} = (s_1, \ldots, s_L)$.

\subsection{Payoff structure}

For $L = 3$ players, the payoff values are given in Table~\ref{tab:payoffs}.
The game is symmetric under cyclic permutation of the players:
\begin{equation}
    h_i(\mathbf{s}) = h_{i+1 \bmod 3}(\sigma(\mathbf{s})),
\end{equation}
where $\sigma$ denotes the corresponding cyclic permutation of the bit string.
Mutual cooperation ($\mathbf{s} = 000$) yields a payoff of 6 per player,
while mutual defection ($\mathbf{s} = 111$) yields only 2---the characteristic tension of the Prisoner's Dilemma.
A single defector receives 10 while the remaining cooperators receive 3 or 0,
creating an incentive to deviate from cooperation.

\begin{table}[h]
    \centering
    \begin{tabular}{cccc}
        \toprule
        Outcome $\mathbf{s}$ & $h_1(\mathbf{s})$ & $h_2(\mathbf{s})$ & $h_3(\mathbf{s})$ \\
        \midrule
        000 & 6  & 6  & 6  \\
        001 & 3  & 3  & 10 \\
        010 & 3  & 10 & 3  \\
        011 & 0  & 6  & 6  \\
        100 & 10 & 3  & 3  \\
        101 & 6  & 0  & 6  \\
        110 & 6  & 6  & 0  \\
        111 & 2  & 2  & 2  \\
        \bottomrule
    \end{tabular}
    \caption{Payoff values for the 3-player quantum Prisoner's Dilemma.}
    \label{tab:payoffs}
\end{table}

For general $L$, players are arranged on a ring with pairwise interactions.
Each player interacts with their two nearest neighbors via a 2-player Prisoner's Dilemma Hamiltonian
$H^{(2)}$ with payoffs $(R, S, T, P) = (3, 0, 5, 1)$,
and the total payoff for player $i$ is the average over incident edges:
\begin{equation}
    H_i = \frac{1}{2}\left( H^{(2)}_{i, i+1} + H^{(2)}_{i, i-1} \right),
\end{equation}
where indices are taken modulo $L$.

\subsection{Matrix product state representation}

For large $L$, storing the full state vector $\ket{\psi} \in (\mathbb{C}^2)^{\otimes L}$ is intractable.
We represent quantum states as matrix product states (MPS) with open boundary conditions:
\begin{equation}
    \ket{\psi} = \sum_{s_1, \ldots, s_L} A_1^{s_1} A_2^{s_2} \cdots A_L^{s_L} \ket{s_1 s_2 \cdots s_L},
\end{equation}
where each $A_j^{s_j}$ is a matrix of dimension $\chi_{j-1} \times \chi_j$ with $\chi_0 = \chi_L = 1$.
The bond dimension $\chi = \max_j \chi_j$ controls the entanglement capacity of the ansatz.
For $\chi = 1$, the state is a product state (zero entanglement);
for $\chi = 2^{\lfloor L/2 \rfloor}$, the MPS can represent any state exactly.

In practice, we work with bond dimension $\chi = 2$, which is sufficient to capture
states with moderate entanglement, including GHZ-type states
$\ket{\mathrm{GHZ}} = (\ket{0}^{\otimes L} + \ket{1}^{\otimes L})/\sqrt{2}$.

\subsection{Nash equilibrium and exploitability}

A state $\ket{\psi}$ constitutes an $\varepsilon$-Nash equilibrium if no player can increase their payoff
by more than $\varepsilon$ through a unilateral unitary deviation:
\begin{equation}
    \max_{U_i \in \mathrm{SU}(2)} \bra{\psi_i(U_i)} H_i \ket{\psi_i(U_i)} - \bra{\psi} H_i \ket{\psi} < \varepsilon \quad \forall\, i,
\end{equation}
where $\ket{\psi_i(U_i)}$ denotes the state with $U_i$ applied only to qubit $i$.
The left-hand side defines the \emph{exploitability} $\varepsilon_i(\psi)$ of player $i$.

We find Nash equilibria using \emph{simultaneous differential best response dynamics}.
At each iteration $n$, we compute the gradient of each player's payoff with respect to
their local unitary:
\begin{equation}
    M_i = \pdv{E_i}{U_i} = \tr_{j \neq i}\left[ H_i \ketbra{\psi}{\psi} \right],
\end{equation}
where the partial trace is over all qubits except $i$.
The unitary update is extracted via the SVD of the matrix $\mathbb{I} - \alpha M_i$:
if $\mathbb{I} - \alpha M_i = Y \Sigma Z^\dagger$, the update is $U_i = Z Y^\dagger$,
which is the unitary closest to the gradient step in operator norm.
All players update simultaneously, and the procedure is iterated until
the total exploitability $\sum_i \varepsilon_i < \varepsilon_{\mathrm{thr}}$.

The exploitability $\varepsilon_i$ is computed by globally optimizing over $\mathrm{SU}(2)$
using differential evolution,
ensuring that the reported Nash equilibria are validated against arbitrary (not just infinitesimal) deviations.

\bibliographystyle{unsrt}
% \bibliography{refs}

\section{Area laws}
\label{sec:area-laws}

A quantum game is defined by the triple \((\hat{h}_i, G_i, N)\) where $N$ is the number of players. We would like to study the quantum advantage associated with the underlying resource state.

\begin{itemize}
    \item In the product state orbit, the game is purely classical (with automatic extension to mixed strategies by controlling the probability amplitudes of individual qubits)
    \item In the volume law cases, we get $\tr (U_i \rho U_i^\dagger \hat{O}) \approx \tr \rho \hat{O}$ for any local observable $\hat{O}$.
    \item So there might be an optimal entanglement structure for the underlying resource state for the maximal quantum advantage (in the sense of mechanism design).
\end{itemize}

\subsection{Introducing non-commutativity}

\begin{itemize}
    \item Product state case: always find a Nash equilibrium irrespective of the payoff Hamiltonians.
    \item Adding entanglement by doing a multi-qubit unitary: can find NE for random Ising game with \(  \hat{h}_L = \hat{h}_0 + V(g_L), \hat{h}_R = \hat{h}_0 + V(g_R) \), where \( V(g) = g \hat{X}_i \) and we randomly pick $g_L, g_R$ from a distribution.
    \item This statement breaks down when we consider the random Heisenberg interactions with \(\hat{h}_i = n_x \hat{X}_i \hat{X}_{i+1} + n_y \hat{Y}_i \hat{Y}_{i+1} + n_z \hat{Z}_i \hat{Z}_{i+1}\). (\(n_x^2 + n_y^2 + n_z^2 = 1\) chosen at random). Adding entanglement in the initial states let the interaction dynamics to end up in a limit cycle.
    \item Note that all games considered so far are translationally invariant, i.e. \(T^\dagger \hat{h}_i T = \hat{h}_i\) 
\end{itemize}



\end{document}
